% ============================================================
\chapter{Th\'eor\`emes de Stabilit\'e}
\label{ch:stabilite}
% ============================================================

L'homologie persistante ne serait qu'une curiosit\'e math\'ematique si elle n'\'etait pas \emph{stable}~: un petit changement dans les donn\'ees doit produire un petit changement dans le diagramme de persistance. Ce r\'esultat fondamental, d\'emontr\'e par David Cohen-Steiner, Herbert Edelsbrunner et John Harer en 2007, est le pilier sur lequel repose toute la TDA appliqu\'ee. Il affirme que la distance bottleneck entre deux diagrammes de persistance est born\'ee par la distance $L^\infty$ entre les fonctions de filtration. Sans ce th\'eor\`eme, le moindre bruit dans les donn\'ees pourrait bouleverser compl\`etement les descripteurs topologiques, les rendant inutilisables en pratique.

\begin{intuition}
Les th\'eor\`emes de stabilit\'e sont le fondement th\'eorique de la
TDA~: ils garantissent que de petites perturbations des donn\'ees
produisent de petits changements dans les diagrammes de persistance.
Sans stabilit\'e, les descripteurs topologiques seraient inutilisables
en pratique.
\end{intuition}

% ============================================================
\section{Distances entre diagrammes~: rappels}
% ============================================================

Avant d'\'enoncer les th\'eor\`emes de stabilit\'e, rappelons les
distances utilis\'ees. Nous les d\'evelopperons au chapitre~\ref{ch:distances}.

\begin{definition}[Distance bottleneck]
Soient $D_1, D_2$ deux diagrammes de persistance. La \emph{distance
bottleneck} est~:
\[
  d_B(D_1, D_2) = \inf_{\gamma} \sup_{x \in D_1}
  \norm{x - \gamma(x)}_\infty
\]
o\`u l'infimum est pris sur toutes les bijections
$\gamma : D_1 \to D_2$ (rappelons que les diagrammes incluent la
diagonale avec multiplicit\'e infinie).
\end{definition}

\begin{definition}[Distance de Wasserstein]
La \emph{distance de Wasserstein d'ordre $q$} ($1 \leq q < \infty$) est~:
\[
  W_q(D_1, D_2) = \left(\inf_{\gamma} \sum_{x \in D_1}
  \norm{x - \gamma(x)}_\infty^q \right)^{1/q}.
\]
\end{definition}

\begin{formules}[title=\textbf{Relation entre distances}]
\[
  d_B(D_1, D_2) = W_\infty(D_1, D_2) \leq W_q(D_1, D_2) \leq
  W_p(D_1, D_2) \quad \text{pour } p \leq q.
\]
\end{formules}

% ============================================================
\section{Th\'eor\`eme de stabilit\'e bottleneck}
% ============================================================

Le th\'eor\`eme fondamental de la TDA est d\^u \`a Cohen-Steiner,
Edelsbrunner et Harer (2007).

\begin{definition}[Fonction de Morse apprivois\'ee]
Une fonction $f : X \to \R$ est dite \emph{apprivois\'ee} (\emph{tame})
si les ensembles de sous-niveaux $f^{-1}((-\infty, r])$ ont une
homologie de dimension finie pour tout $r$, et le nombre de valeurs
critiques est fini.
\end{definition}

\begin{theoreme}[Stabilit\'e bottleneck --- Cohen-Steiner, Edelsbrunner, Harer, 2007]
Soient $f, g : X \to \R$ deux fonctions apprivois\'ees sur un espace
topologique triangulable $X$. Alors~:
\[
  d_B(\mathrm{Dgm}(f), \mathrm{Dgm}(g))
  \leq \norm{f - g}_\infty.
\]
\end{theoreme}

\begin{proof}[Esquisse de preuve]
La preuve repose sur trois ingr\'edients~:
\begin{enumerate}
  \item \textbf{Interpolation}~: consid\'erer la famille
        $f_t = (1-t)f + tg$ pour $t \in [0, 1]$.
  \item \textbf{Lemme de bo\^{\i}te}~: si les diagrammes de $f$ et $g$
        diff\`erent dans une r\'egion, alors $\norm{f - g}_\infty$
        est au moins aussi grande que la taille de cette r\'egion.
  \item \textbf{Construction d'un couplage}~: on construit une bijection
        $\gamma$ entre les diagrammes en suivant les points \`a travers
        l'interpolation.
\end{enumerate}
Voir Cohen-Steiner, Edelsbrunner, Harer (2007) pour la preuve compl\`ete.
\end{proof}

\begin{remarque}
Ce th\'eor\`eme est optimal~: il existe des fonctions pour lesquelles
l'\'egalit\'e est atteinte.
\end{remarque}

% ============================================================
\section{Stabilit\'e pour les nuages de points}
% ============================================================

\begin{corollaire}[Stabilit\'e Vietoris--Rips]
Soient $X, Y \subset \R^d$ deux nuages de points finis. Soit
$d_H(X, Y)$ la distance de Hausdorff. Alors~:
\[
  d_B(\mathrm{Dgm}(\mathrm{VR}(X)), \mathrm{Dgm}(\mathrm{VR}(Y)))
  \leq 2\, d_H(X, Y).
\]
\end{corollaire}

\begin{definition}[Distance de Hausdorff]
La \emph{distance de Hausdorff} entre deux compacts $X, Y$ dans un
espace m\'etrique $(M, d)$ est~:
\[
  d_H(X, Y) = \max\left(
    \sup_{x \in X} \inf_{y \in Y} d(x, y),\;
    \sup_{y \in Y} \inf_{x \in X} d(x, y)
  \right).
\]
\end{definition}

\begin{proof}
Soit $f_X(z) = \inf_{x \in X} d(z, x)$ la fonction distance \`a $X$,
et de m\^eme pour $f_Y$. Alors
$\norm{f_X - f_Y}_\infty = d_H(X, Y)$.
Les diagrammes de persistance des filtrations de sous-niveaux de $f_X$
et $f_Y$ co\"{\i}ncident (aux constantes pr\`es) avec ceux des
filtrations de Rips. Le r\'esultat suit du th\'eor\`eme de stabilit\'e
bottleneck, avec le facteur~$2$ venant de la convention de param\'etrisation.
\end{proof}

% ============================================================
\section{Stabilit\'e de Wasserstein}
% ============================================================

\begin{theoreme}[Stabilit\'e Wasserstein --- Chazal et al., 2009]
Soient $f, g : X \to \R$ deux fonctions apprivois\'ees.
Pour tout $q \geq 1$ et tout $p \in \N$, si les diagrammes $\mathrm{Dgm}_p(f)$
et $\mathrm{Dgm}_p(g)$ ont au plus $N$ points hors de la diagonale, alors~:
\[
  W_q(\mathrm{Dgm}_p(f), \mathrm{Dgm}_p(g))
  \leq C(N, q) \cdot \norm{f - g}_\infty
\]
o\`u $C(N, q) = N^{1/q}$ pour la borne la plus simple.
\end{theoreme}

\begin{remarque}
Contrairement \`a la stabilit\'e bottleneck, la stabilit\'e de Wasserstein
d\'epend du nombre de points dans le diagramme. Des bornes plus fines
existent sous des hypoth\`eses suppl\'ementaires.
\end{remarque}

% ============================================================
\section{Stabilit\'e des repr\'esentations}
% ============================================================

\begin{theoreme}[Stabilit\'e des paysages de persistance]
Les paysages de persistance sont $1$-lipschitziens par rapport \`a la
distance bottleneck~:
\[
  \norm{\lambda_k^{(1)} - \lambda_k^{(2)}}_\infty
  \leq d_B(D_1, D_2)
\]
pour tout $k \geq 1$.
\end{theoreme}

\begin{proposition}[Stabilit\'e des images de persistance]
Les images de persistance sont lipschitziennes par rapport \`a la
distance de Wasserstein $W_1$~: il existe $C > 0$ (d\'ependant de
$\sigma$ et de la fonction de poids) tel que~:
\[
  \norm{\mathrm{PI}(D_1) - \mathrm{PI}(D_2)}_F
  \leq C \cdot W_1(D_1, D_2).
\]
\end{proposition}

% ============================================================
\section{Th\'eor\`eme de stabilit\'e g\'en\'eralis\'e}
% ============================================================

\begin{definition}[Distance d'entrelacement]
Deux modules de persistance $\mathbf{U}$ et $\mathbf{V}$ sont
\emph{$\varepsilon$-entrelac\'es} s'il existe des morphismes naturels
$\varphi : \mathbf{U} \to \mathbf{V}[\varepsilon]$ et
$\psi : \mathbf{V} \to \mathbf{U}[\varepsilon]$ tels que les compos\'ees
$\psi[\varepsilon] \circ \varphi$ et $\varphi[\varepsilon] \circ \psi$
co\"{\i}ncident avec les applications de structure de d\'ecalage
$2\varepsilon$.

La \emph{distance d'entrelacement} est~:
\[
  d_I(\mathbf{U}, \mathbf{V}) = \inf\{\varepsilon \geq 0 :
  \mathbf{U} \text{ et } \mathbf{V} \text{ sont }
  \varepsilon\text{-entrelac\'es}\}.
\]
\end{definition}

\begin{theoreme}[Isom\'etrie alg\'ebrique de stabilit\'e]
Pour tout module de persistance p.f.d.~:
\[
  d_B(\mathrm{Dgm}(\mathbf{U}), \mathrm{Dgm}(\mathbf{V}))
  = d_I(\mathbf{U}, \mathbf{V}).
\]
\end{theoreme}

\begin{remarque}
Ce th\'eor\`eme, d\^u \`a Chazal, Cohen-Steiner, Glisse, Guibas et
Oudot (2009) et Bauer--Lesnick (2015), montre que la distance bottleneck
est la ``bonne'' distance sur les diagrammes de persistance~: elle
co\"{\i}ncide exactement avec la notion naturelle de proximit\'e des
modules de persistance.
\end{remarque}

% ============================================================
\section{Applications de la stabilit\'e}
% ============================================================

\subsection{Inf\'erence topologique}

\begin{theoreme}[Inf\'erence topologique --- Niyogi, Smale, Weinberger, 2008]
Soit $\mathcal{M}$ une sous-vari\'et\'e compacte de $\R^d$ de nombre
d'atteinte $\tau > 0$. Si $X$ est un \'echantillon $\varepsilon$-dense
de $\mathcal{M}$ avec $\varepsilon < \tau / 2$, alors la filtration de
\v{C}ech de $X$ au param\`etre $r \in (\varepsilon, \tau - \varepsilon)$
a le m\^eme type d'homotopie que $\mathcal{M}$.
\end{theoreme}

\subsection{Robustesse au bruit}

\begin{proposition}
Si $X$ est un nuage de points \'echantillonn\'e sur une vari\'et\'e
$\mathcal{M}$ et $X'$ est une version bruit\'ee avec
$d_H(X, X') \leq \delta$, alors~:
\[
  d_B(\mathrm{Dgm}(X), \mathrm{Dgm}(X')) \leq 2\delta.
\]
Les caract\'eristiques topologiques de persistance $> 4\delta$ dans le
diagramme de $X'$ correspondent \`a de vraies caract\'eristiques de
$\mathcal{M}$.
\end{proposition}

% ============================================================
\section{D\'emonstration illustr\'ee}
% ============================================================

\begin{codeblock}[title=\textbf{V\'erification num\'erique de la stabilit\'e}]
\begin{minted}{python}
import numpy as np
from ripser import ripser
from persim import bottleneck

# Cercle exact
theta = np.linspace(0, 2*np.pi, 100, endpoint=False)
X = np.column_stack([np.cos(theta), np.sin(theta)])

# Versions bruitées avec epsilon croissant
epsilons = [0.01, 0.05, 0.1, 0.2, 0.5]
dgm_X = ripser(X, maxdim=1)['dgms']

for eps in epsilons:
    noise = eps * np.random.randn(*X.shape)
    X_noisy = X + noise
    dgm_Y = ripser(X_noisy, maxdim=1)['dgms']

    # Distance bottleneck en H_1
    d_bn = bottleneck(dgm_X[1], dgm_Y[1])

    # Distance de Hausdorff (approximation)
    from scipy.spatial.distance import directed_hausdorff
    d_h = max(directed_hausdorff(X, X_noisy)[0],
              directed_hausdorff(X_noisy, X)[0])

    print(f"eps={eps:.2f}: d_B={d_bn:.4f}, "
          f"d_H={d_h:.4f}, "
          f"ratio={d_bn/d_h:.4f}")
\end{minted}
\end{codeblock}

% ============================================================
\section{Limites de la stabilit\'e}
% ============================================================

\begin{attention}[title=\textbf{Limites}]
\begin{itemize}
  \item La stabilit\'e bottleneck est une borne \emph{sup}~: le
        diagramme peut changer beaucoup moins que la borne.
  \item La borne de Wasserstein d\'epend du nombre de points~: pour
        des diagrammes avec beaucoup de points de faible persistance,
        la borne peut \^etre lache.
  \item La stabilit\'e concerne les \emph{diagrammes}, pas
        n\'ecessairement les \emph{repr\'esentations vectorielles}
        (bien que celles-ci h\'eritent souvent de propri\'et\'es de
        stabilit\'e).
\end{itemize}
\end{attention}

% ============================================================
\section{Exercices}
% ============================================================

\begin{exercice}
Montrer que la distance bottleneck v\'erifie l'in\'egalit\'e
triangulaire.
\end{exercice}

\begin{exercice}
Construire un exemple de deux fonctions $f, g : [0, 1] \to \R$ telles
que $d_B(\mathrm{Dgm}(f), \mathrm{Dgm}(g)) = \norm{f - g}_\infty$
(montrer que la borne est atteinte).
\end{exercice}

\begin{exercice}
V\'erifier num\'eriquement la stabilit\'e bottleneck en g\'en\'erant
des nuages de points sur un tore avec diff\'erents niveaux de bruit.
Tracer $d_B$ en fonction de $d_H$ et v\'erifier la lin\'earit\'e.
\end{exercice}

\begin{exercice}
D\'emontrer que si $\mathbf{U}$ et $\mathbf{V}$ sont
$\varepsilon$-entrelac\'es, alors
$d_B(\mathrm{Dgm}(\mathbf{U}), \mathrm{Dgm}(\mathbf{V})) \leq \varepsilon$.
\end{exercice}

\begin{exercice}
Soit $X$ un nuage de $n$ points i.i.d. sur une vari\'et\'e compacte
$\mathcal{M}$. Montrer que $d_H(X, \mathcal{M}) \to 0$ presque
s\^urement quand $n \to \infty$, et en d\'eduire la convergence des
diagrammes de persistance.
\end{exercice}
